\section{Guenonian Scholasticism}

Julius Evola was among the first writers to take up the ideas of Rene Guenon, and maintained an epistolary relationship with him until the end of Guenon's life. Yet, although Evola wrote a long article in \emph{La Vita Italiana}, titled ``Rene Guenon: A Master for our Times'', in praise of Guenon, he also wrote another article in which he opposed what he called `Guenonian Scholasticism', which he defines as:

\begin{quotex}
This kind of `scholasticism' consists in following passively just about every view ever formulated by Guenon, with a pedantic attitude, without any true investigation or discrimination, and with a real fear to make even the slightest change in the master's formulations.
\end{quotex}


This accusation is directed toward Guenon's followers as much as against the Master himself. We can agree with Evola that Guenon's writings are not revealed text, but we also need to take into account the Guenon offered solid reasons for his perspectives. Therefore, we should expect the same from Evola, and not assume that Guenon is writing from his ``personal equation''. In our day, ironically, we need to oppose an Evolian Scholasticism that takes some of the more marginal positions of Evola as representative of the whole of Tradition. What follows are the points espoused by Guenon that Evola takes issue with.


\paragraph{Knowledge and Action}


Evola disputes Guenon's contention that contemplation, or knowledge, is superior to action. Guenon explains his position most fully in \emph{The Crisis of the Modern World}, where he associates knowledge with the Brahmin caste and the East, while associating action with the Kshatriyas and the West. If it were just a question of two perspectives, that would be one thing, but Guenon goes on to infer that the West is deviant in that regard. Evola takes the first position (that there are two legitimate paths, one knowledge and one of action), but then in the next paragraph contradicts himself by giving action primacy.

Guenon's contention is that the principle of Action lies outside Action, viz., in the Unmoved Mover, a position also held by Evola. Guenon asserts that the Unmoved Mover is actually knowledge. Evola doesn't actually address Guenon's argument, so we'll give it a try. The debate is a little puzzling, if only because Guenon asserts that there is a more fundamental principle that unites the Brahmin (knowledge) and the Kshatriya (action), but then seems to have forgotten it. Both Guenon and Evola appeal to the Bhagavad Gita as their authority and come to different conclusions. On the other hand Tilak\footnote{\url{http://www.gornahoor.net/?p=1619}}, who was respected by both men, seems to have learned the right lesson. It is the noumenal Will that is the Unmoved Mover and the uniting principle. While it is true that there have been different paths, Boris Mouravieff describes a Fourth Way path, based on Hermetic and Hesychast Traditions in the Eastern church, that combine both. \emph{Pace} Evola, knowledge is essential, but \emph{pace} Guenon, knowledge without a corresponding change in the level of being is ineffective. Acquiring Knowledge is the easier part, while the change in Being is more difficult as it involves the Will. Evola mentions the practical and ascetic attitude of Early Buddhism to bolster his point.

\paragraph{Spiritual Authority}


Evola challenges Guenon's position of Spiritual Authority vis-a-vis Regal Authority. Again, this seems like a non-issue, because it also depends on that deeper principle uniting the two castes and seems to not to take into account Temporal Power. Yes, the priestly caste ought not to exercise temporal power over the administrative caste, yet they cannot abandon their teaching task on spiritual matters. Nevertheless, Guenon and Evola agree that there is a Regal initiation that integrates the two castes. For example, it was accepted that the Pope should have temporal power over the Papal States, but I don't recall Evola ever mentioning it.

I think the dispute is exacerbated by historical circumstances, in which the Catholic religion came to be dominated by theological disputes which overflowed into the secular realm. This may be because Christianity did not have a Law, \emph{per se}, as do Traditional religions. Contrast this with the Roman idea of the Emperor or the Muslim idea of the Caliphate. The Emperor or Caliphate is supposed to represent, not just the political unity, but also the spiritual unity of the State. However, in these cases, this involves ensuring that the religious law is followed and enforced. Theological and philosophical positions are personal matters in those systems, so the enforcement of creedal orthodoxy is not their task.

\paragraph{Initiatory Organizations and Masonry}


Guenon was overly punctilious about the regularity of initiatory organizations and he accepted Masonry as one such. Here we agree with Evola. Self-initiation is possible as we have pointed out\footnote{\url{http://www.gornahoor.net/?p=1721}}. Also, even if the initiatory impetus of a Tradition is passive and virtual, it can still be made active. Evola seems to have picked up his strong \emph{animus} against Masonry from the long-time Catholic opposition to it ... part of \emph{his} ``personal equation.'' This has two roots. The first is related to the Templars which was unjustly destroyed. This led to the end of initiatory organizations in the West, even if Masonry carried it on for a time. Nevertheless, there was a time when Masonry did indeed transform itself into a revolutionary and subversive organization. Guenon describes that process in some detail.

\paragraph{Buddhism}


Guenon devalued Buddhism, but later corrected himself by discovering an alleged ``Brahmanic'' version of Buddhism. This Evola rejects; it is sufficient to refer to his work \emph{The Doctrine of Awakening}. In Evola's conception, initiation on this path depends on the descent of ``forces from above'', something we are sympathetic with for obvious reasons.

\paragraph{The Necessity of a Traditional Exoterism}


They both make sense. From Guenon's perspective, in a Traditional society, one simply accepts the exoteric system of the group (this is what led to the conviction of Socrates). However, Evola rightly points out that the gap in our day between Esoterism and the available Exoterisms is currently unbridgeable. Perhaps the way forward is to find a way to bridge them.


\flrightit{Posted on 2011-03-03 by Cologero}

\begin{center}* * *\end{center}

\begin{footnotesize}
\begin{sffamily}
\texttt{HOO on 2011-03-03 at 05:20 said:}

The Brahmin vs. Kshatriya discussion is most of the time ridiculous and useless, and entertaining it seems like a symptom of counter-productive bookishness, for just as the Kshatriya had have suffered a gradual reduction to ``temporal power'', to administrative, juristic and military functions, the Brahmin had suffered a gradual reduction to ``wordy or moral temporal pressure'', to administrative, juristic and psuedoritualistic functions (to which Buddhism was in a way a response to).

´By far more than the sacerdotal title of `saint', it is that of Lord which, of all times, was used by all people to designate `God' -- as symbol of the highest metaphysical state in which the human being can integrate. Also Guénon, when he uses the title `King of the World' to designate the supreme centre of spiritual authority, is he not himself referring to a non-sacerdotal dignity? Does he not himself note (p. 137) the symbolic relation between the sceptre, emblem of royal dignity, and the `axis of the world'?´ Evola, ``Spiritual Authority and Temporal Power'' Many Guénonianists seem horrified by the truth that even the most vile serial killer (e.g. Angulimala), criminal or cannibal can become enlightened. They seem to be somehow attached to the mere-human cycles of moralism and rationalism, or such social or geometrical judgment.

´Therefore, a good Christian should beware that mathematicians, and any others who prophesy impiously... may be entangled in the companionship of demons*.´ -- St. Augustine

*I'll illuminate some of you to the truth that a demon can mean a perspective or an angle (and the inescapable prejudice of an idea).

If the reader finds he react in some way emotionally to this comment he has found at least one of his demons.

\hfill

\texttt{Cologero on 2011-03-03 at 08:32 said:}

Yes, moralism and sentimentalism, things Guenon seeks to avoid, are the result when the ideal of contemplation is reduced to the human level. I think another problem is that in the West, spiritual progress is often considered equivalent to holding the correct beliefs or opinions.

That is why I recommend the book ``Opening the Dragon Gate'', since Taoism was so influential on both Guenon and Evola. The young Taoist initiate reveals how he had to undergo several trials before being taught anything that we would consider ``knowledge''. This is Evola's point about early Buddhism, which, by the way, interacted with Taoism. But here again is the point. Tao means ``way'' or ``path'', Buddhism has an 8-fold path. The point is to follow the path, not to believe certain things. That is why there could be eight schools of Buddhism in China without a religious war ensuing.

\hfill

\texttt{EXIT on 2011-03-03 at 09:21 said:}

Far from useless the dispute is critical for the reduction to a feminine spirituality, i.e.., moralism, sentimentalism, and empty ritualism, occurs as Guenon points out at the moment when an inversion of hierachy takes place due to the revolt by the regal warrior caste.

\hfill

\texttt{EXIT on 2011-03-03 at 09:28 said:}

Hoo, you are wrong to say that Guenonians are horrified by immorality. But I don't know any serial killers or cannibals that are truly enlightened, do you?

I would also take issue with Christian polemics, as if Augustine was an accurate judge of Greek, Pythagorean, and pagan traditions. I couldn't care less what this or that desert father said, and I am sure that they didn't believe much of their own arguments regarding such.

\hfill

\texttt{kadambari on 2011-03-03 at 10:13 said:}

@HOO

´Therefore, a good Christian should beware that mathematicians, and any others who prophesy impiously... may be entangled in the companionship of demons*.´ -- St. Augustin

I am unable to understand the above comment? Why does he think Mathematicians prophesy impiously? As EXIT pointed out, how much does he really understand of Pythagoras, and the other pre-Christians? Is this not a Christian prejudice towards science? I am understanding science in the philosophical sense as the ancients who created and contributed to science understood it, as a part of the contemplative philosophical life... \hfill

\texttt{Ken Hoop on 2011-03-03 at 15:40 said:}

\url{http://zennist.typepad.com/zenfiles/2011/03/buddhism-is-reformed-brahmanism.html}


Evola would not have agreed with this?

\hfill

\texttt{Will on 2011-03-04 at 11:34 said:}

Kadambari, I think that what Augustine means by ``mathematicians'' is something more along the lines of ``numerologists.'' I don't know how much he did or didn't understand Pythagoreanism and Paganism, but I don't doubt that there were indeed people in his time who claimed to be prophets but were charlatans. One thing early Christianity has in common with early Buddhism is a disdain for divination. The Buddha said about astrology, ``What can the stars do to one who is awake?''

Incidentally, there's an amusing scene in the movie Pi, in which two mathematicians are talking and one of them has become interested in Kabbalist mathematics. The other mathematician warns him that if he goes down that road, ``You're no longer a mathematician -- you're a numerologist,'' and he enunciates ``numerologist'' with venomous contempt.

\hfill

\texttt{EXIT on 2011-03-04 at 14:05 said:}

HOO is implying that a traditional application of principles by Guenon is a rigid mathematical rationalism which doesn't hold up to reality. But Reality is the ultimate Mathematician regardless what the idol-fearing Christians, who are frightened that the devil may be hiding under every rock, think.

\hfill

\texttt{Martin on 2019-05-10 at 03:48 said:}

What ``Brahmanic'' Buddhism did Guenon claim to find? Can you kindly point me to a source where he does this? Thank you.

\hfill

\texttt{Cologero on 2019-05-10 at 04:29 said:}

Evola makes that claim in ``René Guénon: a Teacher for Modern Times''

The topic is also discussed in ``The Essential Rene Guenon''. Guenon claimed that the original Buddhism of India is different from what it eventually became.

\hfill

\end{sffamily}
\end{footnotesize}

